\chapter{sysfs 与 GPIO 按键驱动}
\label{chap:sysfs}
本章按当前代码说明两件新加的东西：\file{/sys} 怎样把内存里的 driver core（kobject、device、bus、driver）投影成目录和只读文件；100ASK T113-S3 上 RESET 旁的 USER KEY怎样经 \code{gpio-keys} 驱动变成 \file{/dev/key0} 的按键事件，以及它在 \file{/sys} 里的属性。

\begin{table}[H]
\centering
\small
\begin{tabularx}{\linewidth}{>{\raggedright\arraybackslash}X>{\raggedright\arraybackslash}X}
\toprule
文件 & 内容 \\
\midrule
\code{src/kernel/fs/sysfs.c} & sysfs 文件系统：lookup、readdir、read \\
\code{src/include/fs/sysfs.h} & \code{t113\_sysfs\_init()} \\
\code{src/fs/Kconfig} & \code{CONFIG\_SYSFS}，T113 板默认打开 \\
\code{src/include/t113/device.h} & 新增 \code{struct device\_attribute}，\code{struct device} 加 \code{attributes/num\_attributes} \\
\code{src/boards/100ask\_t113/device.c} & 解绑或 probe 失败时清空属性表 \\
\code{src/drivers/input/gpio\_keys.c} & 按键驱动 \\
\code{src/include/input/gpio\_keys.h} & 平台数据和 \code{struct minios\_input\_event} \\
\code{src/boards/100ask\_t113/board.c} & \code{user-key} 平台设备，PB4 \\
\code{src/boards/100ask\_t113/gpio.c} & GPIO 外部中断寄存器操作 \\
\code{user-t113/keytest.c}、\code{user-t113/include/input.h} & 用户态测试程序和事件结构 \\
\code{tools/mkfs\_minifs.c} & 在根文件系统里建空目录 \file{/sys} 作为挂载点 \\
\bottomrule
\end{tabularx}
\end{table}

\section{sysfs}

\subsection{挂载时机}

\begin{txtcode}[basicstyle=\ttfamily\scriptsize]
kernel_boot()                                    kernel/init/main.c
+-- t113_minifs_mount_rootdev()      /
+-- t113_fat32_mount_vfs()           /boot
+-- t113_chrdev_init()
+-- t113_devfs_init()                /dev
+-- ...
+-- t113_driver_core_init()          建 devices/bus/drivers 三个 kset，注册设备和驱动，扫 MMC
+-- t113_sysfs_init()                /sys，失败 → PANIC "cannot mount sysfs at /sys"
    +-- t113_vfs_mount_fs("/sys", 0x62656572, 0, 0, &sysfs_iops, &sysfs_dir_fops)
\end{txtcode}

\code{0x62656572} 是 Linux 的 \code{SYSFS\_MAGIC}（ASCII "beer"）。sysfs 不保存任何数据：\code{super\_block} 没有 \code{s\_fs\_info}，也没有 \code{super\_operations}。每次 lookup、readdir、read 都直接去查 driver core 当时的状态，所以先挂后挂都可以；放在\code{t113\_driver\_core\_init()} 之后只是为了让挂载时树已经完整。

mount 时 VFS 要先在 minifs 里找到 \file{/sys} 目录，所以 \code{mkfs\_minifs} 新建了一个空的\code{sys} 目录。旧的 \code{minifs.img} 没有它，必须重新 \code{make t113} 并烧写分区 2，否则启动会停在上面的 PANIC。

\subsection{目录结构}\label{sec:sysfs-tree}

sysfs 只认两棵子树：

\begin{txtcode}
/sys
+-- devices/                         按 device.parent 展开的设备拓扑
|   +-- <dev>/[<child>/...]
|       +-- name   bus   driver   path       每个设备都有的 4 个属性
|       +-- <驱动提供的属性>                    device->attributes[]
+-- bus/
    +-- <bus>/
        +-- name
        +-- devices/<dev>/                   与 /sys/devices 下同一设备的属性相同
        +-- drivers/<drv>/
            +-- name
            +-- bound                        绑定到这个驱动的设备名，空格分隔
\end{txtcode}

下面是主机侧测试的实际输出（见第~\ref{sec:sysfs-test} 节；注册了 \code{uart3}、\code{timer0}、\code{user-key} 三个平台设备和 \code{t113-uart}、\code{gpio-keys} 两个驱动）。板子上还会有 \code{gic}、\code{twi2}、\code{sdc1}、\code{csi0}，以及 \code{i2c}、\code{sdio} 两条总线和它们下面的设备。

\begin{txtcode}
/sys/
  devices/
    platform/
      uart3/
        name  = "uart3"
        bus  = "platform"
        driver  = "t113-uart"
        path  = "/devices/platform/uart3"
      timer0/
        name  = "timer0"
        bus  = "platform"
        driver  = "(none)"
        path  = "/devices/platform/timer0"
      user-key/
        name  = "user-key"
        bus  = "platform"
        driver  = "gpio-keys"
        path  = "/devices/platform/user-key"
        value  = "0"
        keycode  = "28"
      name  = "platform"
      bus  = "(none)"
      driver  = "(none)"
      path  = "/devices/platform"
  bus/
    platform/
      devices/
        uart3/ ...  timer0/ ...  user-key/ ...      内容同上
      drivers/
        t113-uart/
          name  = "t113-uart"
          bound  = "uart3"
        gpio-keys/
          name  = "gpio-keys"
          bound  = "user-key"
      name  = "platform"
\end{txtcode}

\file{/sys/devices/platform} 自己也是一个 \code{struct device}（\code{platform\_bus}），所以它下面既有子设备目录，也有 4 个属性文件。

\subsection{路径怎样对应到 driver core}

kobject 的路径由 \code{kobject\_get\_path()} 沿 \code{parent} 链拼出。\code{device\_add()} 把每个设备加进 \code{devices} kset；没有 parent 的 kobject 以 kset 自己的 kobject 为 parent，所以平台设备的路径是 \file{/devices/platform/<name>}（\code{platform\_device\_register()} 把\code{dev.parent} 设成 \code{platform\_bus}）。

sysfs 把 \file{/sys/devices/a/b} 去掉 \file{/sys} 后得到 \file{/devices/a/b}，再在 \code{devices} kset里找路径相同的 kobject：

\begin{txtcode}
split_path("/sys/devices/platform/user-key/value")
  → part[] = { "devices", "platform", "user-key", "value" }，最多 8 级，每级 ≤ 30 字节
sysfs_path_type()
  +-- device_model_path(全部 4 级) = "/devices/platform/user-key/value"
  |   find_device_path() → 没有这个设备
  +-- device_model_path(前 3 级)   = "/devices/platform/user-key"
  |   find_device_path() → user-key 的 struct device（get_device 加引用）
  +-- 最后一级 "value"：不是 name/bus/driver/path，查 device->attributes[] → 找到
  +-- 返回 VFS_NODE_FILE，put_device()
\end{txtcode}

\file{/sys/bus/...} 则按层数解析：

\begin{table}[H]
\centering
\small
\begin{tabularx}{\linewidth}{>{\raggedright\arraybackslash}Xl>{\raggedright\arraybackslash}X}
\toprule
路径 & 类型 & 查法 \\
\midrule
\file{/sys/bus} & 目录 & 固定 \\
\file{/sys/bus/<bus>} & 目录 & \code{kset\_find\_obj(buses, name)} \\
\file{/sys/bus/<bus>/name} & 文件 & \code{bus->name} \\
\file{/sys/bus/<bus>/devices}、\code{drivers} & 目录 & 固定 \\
\file{/sys/bus/<bus>/devices/<dev>} & 目录 & \code{bus\_find\_device\_by\_name()} \\
\file{/sys/bus/<bus>/devices/<dev>/<attr>} & 文件 & 同上，再查属性名 \\
\file{/sys/bus/<bus>/drivers/<drv>} & 目录 & \code{bus\_for\_each\_drv()} 按名字找 \\
\file{/sys/bus/<bus>/drivers/<drv>/name}、\code{bound} & 文件 & 同上 \\
\bottomrule
\end{tabularx}
\end{table}

\subsection{三个 VFS 操作}

sysfs 只有一张 \code{inode\_operations} 和两张 \code{file\_operations}：

\begin{ccode}
static const struct inode_operations sysfs_iops = { .lookup = sysfs_lookup };
static const struct file_operations sysfs_dir_fops = { .readdir = sysfs_readdir };
static const struct file_operations sysfs_file_fops = { .read = sysfs_read };
\end{ccode}

\textbf{lookup。} \code{sysfs\_lookup(dir, child)} 用 \code{child->d\_path} 跑一遍\code{sysfs\_path\_type()}，存在就 \code{t113\_vfs\_bind\_inode()}：\code{i\_ino} 是整条路径的 32 位FNV-1a 哈希，\code{i\_size = 0}，目录配 \code{sysfs\_dir\_fops}，文件配 \code{sysfs\_file\_fops}。

\textbf{readdir。} \code{sysfs\_readdir(inode, path, index, entry)} 现算第 \code{index} 项：

\begin{txtcode}
/sys                         index 0 = devices，1 = bus
/sys/devices[/...]           先列 parent 等于当前设备的子设备（kset_for_each），
                             再列 4 个公共属性和 device->attributes[]
/sys/bus                     buses kset 里的第 index 个 bus
/sys/bus/<bus>               devices、drivers、name
/sys/bus/<bus>/devices       bus_for_each_dev() 的第 index 个设备
/sys/bus/<bus>/drivers       bus_for_each_drv() 的第 index 个驱动
/sys/bus/<bus>/devices/<d>   4 个公共属性 + 驱动属性
/sys/bus/<bus>/drivers/<d>   name、bound
\end{txtcode}

\textbf{read。} \code{sysfs\_read(file, buf, n)} 每次调用都按 \code{file->f\_dentry->d\_path} 重新生成整个值（最多 256 字节），再从 \code{f\_pos} 处拷贝。读到末尾返回 0，所以\code{cat} 一类的循环读能正常结束。

\begin{txtcode}
sysfs_read()
+-- sysfs_generate(path)
    +-- /sys/devices/.../<attr>              show_device(device, attr)
    +-- /sys/bus/<bus>/name                  bus->name + "\n"
    +-- /sys/bus/<bus>/devices/<d>/<attr>    show_device()
    +-- /sys/bus/<bus>/drivers/<d>/<attr>    show_driver()：name，或 bound 列表
show_device()
+-- name / bus / driver                      字符串 + "\n"，没有时写 "(none)"
+-- path                                     kobject_get_path() + "\n"
+-- 其他                                     device->attributes[i].show(dev, buf, cap)
\end{txtcode}

sysfs 没有 \code{write}，\code{t113\_vfs\_open()} 对 \code{O\_RDWR} 直接拒绝，所以 \file{/sys} 是只读的。

\subsection{驱动怎样加属性}

\code{device.h} 新增：

\begin{ccode}
struct device_attribute {
    const char *name;
    int (*show)(struct device *dev, char *buffer, uint32_t capacity);
};

struct device {
    ...
    const struct device_attribute *attributes;   /* 驱动 probe 时设置 */
    uint32_t num_attributes;
};
\end{ccode}

驱动在 probe 成功前把静态数组挂到 \code{dev->attributes}。\code{show} 返回写入的字节数，出错返回负数。driver core 在三处把它清零：\code{device\_initialize()}、probe 失败的\code{really\_probe()}、解绑的 \code{device\_release\_driver\_locked()}。属性的生命周期因此和“设备绑定了这个驱动”一致，解绑后 \file{/sys} 里不会留下指向已释放驱动数据的文件。

\subsection{引用和锁}

\begin{itemize}
  \item lookup 和 readdir 在 VFS 的 \code{vfs\_lock} 里运行（由 \code{lookup\_path()} 和\code{t113\_vfs\_readdir()} 调用）；\code{read} 不在这把锁里，\code{t113\_vfs\_read()} 调驱动前已经放开了它。
  \item 找到的 \code{device}、\code{bus\_type}、\code{device\_driver} 都先加 kobject 引用（\code{get\_device}、\code{kset\_find\_obj}、\code{kobject\_get}），用完马上放掉，不在 inode 里长期持有。所以设备删除后，下一次访问就会看到它不存在了。
  \item 遍历时持有的锁是 kset 的 spinlock 或 driver core 的 \code{core\_lock}，回调里只拿kobject 自己的 \code{ref.lock}，没有同一把锁重入。
  \item 驱动的 \code{show} 由 \code{sysfs\_read()} 调用，在 \code{vfs\_lock} 之外。\code{readdir} 不调用\code{show}。\code{gpio\_keys} 的属性只做单次 GPIO/计数器读取，不持有 VFS 全局锁。
\end{itemize}

\subsection{与 Linux sysfs 的差别}

\begin{table}[H]
\centering
\small
\begin{tabularx}{\linewidth}{>{\raggedright\arraybackslash}X>{\raggedright\arraybackslash}X>{\raggedright\arraybackslash}X}
\toprule
方面 & Linux & miniOS \\
\midrule
实现 & kernfs：每个 kobject、属性在注册时建 \code{kernfs\_node} & 不建节点，按路径实时查 driver core \\
设备属性 & \code{DEVICE\_ATTR\_RO()} + \code{sysfs\_create\_group()}，可读写 & \code{dev->attributes[]}，只读 \\
\code{driver}、\code{subsystem} & 符号链接 & \code{driver}、\code{bus} 是文本章件 \\
\file{/sys/bus/*/devices/*} & 指向 \file{/sys/devices/...} 的符号链接 & 同样内容的目录，\code{path} 属性写真实路径 \\
驱动目录 & 每个设备一个符号链接，加 \code{bind}、\code{unbind}、\code{uevent} & \code{name} 和 \code{bound} 两个文本章件 \\
\file{/sys/class}、\file{/sys/module}、\code{uevent} & 有 & 没有 \\
inode 号 & kernfs 分配，唯一 & 路径的 32 位哈希 \\
\bottomrule
\end{tabularx}
\end{table}

\section{GPIO 按键驱动}

按键数据从硬件到用户态只走一条链路：

\begin{txtcode}
USER KEY（PB4，低电平有效）
  → PB4 EINT pending bit 4
  → GPIOB_NS（GIC IRQ 101）
  → IRQ 异常入口和通用 IRQ 分发表
  → gpio_keys_irq() 写入 input_event 环形队列
  → 唤醒阻塞读取 /dev/key0 的进程
  → keytest 打印 pressed/released
\end{txtcode}

\subsection{硬件和平台数据}

USER KEY 接在 PB4，按下接地，平时靠内部上拉保持高电平。\code{board.c} 在\code{CONFIG\_INPUT\_GPIO\_KEYS=y} 时注册平台设备：

\begin{ccode}
static const struct resource user_key_resources[] = {
    DEFINE_RES_IRQ(T113_IRQ_GPIOB),             /* 101 = GIC SPI 69 */
};
static const struct gpio_keys_platform_data user_key_pdata = {
    .gpio = T113_GPIO(T113_PB, 4),              /* 1 * 32 + 4 = 36 */
    .irq = T113_IRQ_GPIOB,
    .code = MINIOS_KEY_ENTER,                   /* 28，与 Linux KEY_ENTER 相同 */
    .active_low = 1u,
    .pull = T113_PULL_UP,
    .debounce_ms = 20u,
};
static struct platform_device user_key_device = {
    .name = "user-key",
    .compatible = "gpio-keys",
    ...
};
\end{ccode}

\code{gpio\_keys\_driver} 用 \code{of\_match\_table = \{ "gpio-keys" \}} 与它匹配，驱动名是\code{gpio-keys}。

用到的 PIO 寄存器（\code{T113\_PIO\_BASE = 0x02000000}，每个 GPIO 端口 0x30 字节，外部中断区从 0x200 开始，每个硬件端口号占 0x20 字节）。T113 没有实现普通 GPIO 端口 PA，但硬件仍保留 bank 0；PB 的硬件端口号是 1，所以 PB EINT 从 \code{0x200 + 1 * 0x20 = 0x220} 开始，不能把 PB 压缩成 bank 0。

\begin{table}[H]
\centering
\small
\begin{tabularx}{\linewidth}{lll>{\raggedright\arraybackslash}X}
\toprule
寄存器 & 地址 & PB4 的位 & 写入值 \\
\midrule
\code{PB\_CFG0} & 0x02000030 & [19:16] & 0xE = \code{T113\_GPIO\_IRQ}（PB\_EINT4） \\
\code{PB\_DAT} & 0x02000040 & [4] & 读电平 \\
\code{PB\_PULL0} & 0x02000054 & [9:8] & 1 = 上拉 \\
\code{PB\_EINT\_CFG0} & 0x02000220 & [19:16] & 4 = 双边沿 \\
\code{PB\_EINT\_CTL} & 0x02000230 & [4] & 1 = 使能 \\
\code{PB\_EINT\_STATUS} & 0x02000234 & [4] & 读 pending，写 1 清除 \\
\code{PB\_EINT\_DEB} & 0x02000238 & 整个端口 & 0x71（24 MHz、/128） \\
\bottomrule
\end{tabularx}
\end{table}

实板调通暴露了两个互相独立的问题：最初把旧款 sunxi 常见的 EINT mux 6 沿用到了 T113，但 T113/D1 应为 0xE；同时最初用 \code{(port - PB) * 0x20} 计算 EINT bank，把 PB 错放进保留的 PA 槽 0x200。两处都修正后 PB4 才能真正置位 pending。

\subsection{probe}

\begin{txtcode}
driver_register(gpio-keys) → platform 总线 match → really_probe()
+-- gpio_keys_probe(pdev)
    +-- pdata = dev.platform_data；platform_get_irq(pdev, 0) 必须等于 pdata->irq
    +-- key0 = { head = tail = 0, last_irq_us = 0 }，初始化 spinlock 和等待队列
    +-- t113_gpio_set_pull(PB, 4, 上拉)
    +-- t113_gpio_direction_input(PB4)；读一次电平作为初始 value
    +-- t113_chrdev_register(13, 0, &gpio_keys_fops)       设备号 → 驱动 fops
    +-- t113_devfs_create_chrdev("key0", 13, 0)             /dev/key0 → 设备号
    +-- t113_request_irq(101, gpio_keys_irq, &key0)
    +-- t113_gpio_irq_config(PB4, BOTH)
    |   +-- 关 PB4 中断，写 EINT_CFG0，写 EINT_DEB = 0x71
    |   +-- PB_CFG0 切到 EINT 功能
    |   +-- 清掉 pending
    +-- platform_set_drvdata(pdev, &key0)
    +-- dev.attributes = { value, keycode, irq_count, pending, mux }
    +-- t113_gpio_irq_enable(PB4)；t113_gic_enable_irq(101)
    +-- "user-key: PB4 active-low key -> /dev/key0 (EINT mux=0x0e)"
\end{txtcode}

major 13 沿用 Linux 输入子系统的主设备号（\code{INPUT\_MAJOR}）。驱动状态是一个静态\code{key0}，所以现在只支持一个按键。

\subsection{中断路径}

\begin{txtcode}
PB4 电平变化
+-- PB_EINT_STATUS[4] = 1 → GPIOB_NS → GIC SPI 69 → IRQ 101
    +-- kernel_irq_entry 保存现场 → t113_irq_handler() → t113_irq_dispatch(101)
    +-- gpio_keys_irq(irq, &key0)
        +-- PB_EINT_STATUS 第 4 位没置 → 返回（GPIOB 整组共用这一个中断）
        +-- 写 1 清 pending
        +-- now = t113_time_us()；读电平；active_low 取反得到 value
        +-- value 与上次相同，或距上次被接受的事件不到 20 ms → 丢弃
        +-- spin_lock_irqsave(&key0.lock)
        |   +-- last_irq_us = now；key0.value = value
        |   +-- 队列没满：events[head] = { now/1000, EV_KEY, 28, value }，head++
        +-- spin_unlock_irqrestore
        +-- wake_up_all(&key0.wait)
\end{txtcode}

事件队列是 16 项的环形缓冲区，\code{head == tail} 为空，\code{head + 1 == tail} 为满，满了新事件直接丢掉（实际能存 15 个）。

\subsection{阻塞读}

\begin{txtcode}
read(fd, &ev, sizeof ev)                    keytest
+-- sys_read → t113_vfs_read(file)
    +-- file->f_op->read = gpio_keys_read  （chrdev_open 已把 f_op 换成驱动 fops）
        +-- count < 12 → -1
        +-- 循环：
            +-- spin_lock_irqsave(&key0.wait.lock)；spin_lock(&key0.lock)
            +-- 队列非空 → 取出一个事件，解锁，拷给用户，返回 12
            +-- 队列空   → 解开 key0.lock，wait_queue_sleep_locked()
                           （入等待队列、解开 wait.lock、t113_sched_yield()）
\end{txtcode}

先拿等待队列的锁、关着中断检查队列，再入队睡眠，所以“检查为空”和“开始睡眠”之间中断插不进来，不会丢唤醒。字符设备不走 \code{f\_pos\_lock}，所以一个进程阻塞在\file{/dev/key0} 上不会挡住别的 fd。一次 \code{read} 只返回一个事件。

事件结构（内核 \code{struct minios\_input\_event} 与用户态 \code{struct input\_event} 布局相同，12 字节）：

\begin{ccode}
struct input_event {
    uint32_t time_ms;     /* t113_time_us() / 1000 */
    uint16_t type;        /* EV_KEY = 1 */
    uint16_t code;        /* KEY_ENTER = 28 */
    int32_t value;        /* 1 按下，0 松开 */
};
\end{ccode}

\subsection{sysfs 属性}

\begin{table}[H]
\centering
\small
\begin{tabularx}{\linewidth}{>{\raggedright\arraybackslash}X>{\raggedright\arraybackslash}X>{\raggedright\arraybackslash}X}
\toprule
文件 & \code{show} & 内容 \\
\midrule
\file{/sys/devices/platform/user-key/value} & \code{gpio\_keys\_show\_value()} & 直接读取 PB4：0 松开、1 按下 \\
\file{/sys/devices/platform/user-key/keycode} & \code{gpio\_keys\_show\_keycode()} & 28 \\
\file{/sys/devices/platform/user-key/irq\_count} & \code{gpio\_keys\_show\_irq\_count()} & 已进入 PB4 handler 的原始中断次数 \\
\file{/sys/devices/platform/user-key/pending} & \code{gpio\_keys\_show\_pending()} & PB\_EINT\_STATUS[4] 当前值 \\
\file{/sys/devices/platform/user-key/mux} & \code{gpio\_keys\_show\_mux()} & PB4 mux，正常为十进制 14（0xE） \\
\bottomrule
\end{tabularx}
\end{table}

同样的五个文件也出现在 \file{/sys/bus/platform/devices/user-key/} 下。

\begin{txtcode}
/sys/devices/platform/user-key   看设备、总线、绑定的驱动和当前状态
/dev/key0                        阻塞读按键事件
\end{txtcode}

\subsection{用户态}

下面的输出是示例，时间值随实际按键变化：

\begin{shcode}
$ keytest
keytest: press the USER key beside RESET
key code=28 pressed time=14811 ms
key code=28 released time=14994 ms
$ cat /sys/devices/platform/user-key/value
0
$ cat /sys/devices/platform/user-key/mux
14
$ cat /sys/devices/platform/user-key/irq_count
2
$ cat /sys/bus/platform/drivers/gpio-keys/bound
user-key
\end{shcode}

\subsection{与 Linux gpio-keys 的差别}

\begin{table}[H]
\centering
\small
\begin{tabularx}{\linewidth}{l>{\raggedright\arraybackslash}X>{\raggedright\arraybackslash}X}
\toprule
方面 & Linux & miniOS \\
\midrule
设备描述 & 设备树 \code{gpio-keys} 节点，每个子节点一个键 & \code{board.c} 里的平台数据，只有一个键 \\
设备文件 & 输入子系统分配 \file{/dev/input/eventN}（major 13） & 固定 \file{/dev/key0}（13:0） \\
事件 & \code{struct input\_event \{ timeval, type, code, value \}}，每次变化后跟 \code{EV\_SYN} & 12 字节，毫秒时间，没有 \code{EV\_SYN} \\
消抖 & \code{gpiod\_set\_debounce()}，否则延迟工作队列在消抖时间后重新采样 & 在中断里比较时间戳，窗口内的边沿直接丢 \\
读 & 一次可读多个事件，支持 \code{O\_NONBLOCK}、\code{poll} & 一次一个，只能阻塞 \\
GPIO IRQ 分发 & GPIO irqchip 把 bank 父 IRQ 解复用为每个 pin 的虚拟 IRQ & gpio-keys 直接占用 GPIOB 父 IRQ 101 \\
\bottomrule
\end{tabularx}
\end{table}

\section{测试和发现的问题}\label{sec:sysfs-test}

\code{make t113} 在 \code{-Wall -Wextra} 下没有警告。另外把 \code{vfs.c}、\code{sysfs.c}、\code{kobject.c}、\code{device.c}、\code{platform\_bus.c}、\code{gpio\_keys.c} 用主机 gcc（ASan/UBSan）编在一起，GPIO、IRQ、时钟用桩函数代替，按板子的顺序注册设备和驱动后：

\begin{itemize}
  \item 递归 \code{readdir} + \code{read} 整个 \file{/sys}，得到第~\ref{sec:sysfs-tree} 节的树，每个属性都能读；
  \item \file{/sys/nope}、\file{/sys/devices/platform/uart3/value} 打不开，\code{O\_RDWR} 打开属性被拒；
  \item 结束时 VFS 只剩开机时的 3 个 dentry，没有遗留的 file；
  \item 模拟按键：拉低 PB4 并触发中断后 \code{value} 从 0 变成 1。
  \item 100ASK T113-S3 实板：PB4 按下和释放分别得到 code 28、value 1/0，时间戳正常。
\end{itemize}

发现的问题：

\begin{enumerate}
  \item \textbf{消抖会丢掉最后一个边沿，按键事件状态可能卡在“按下”。} 中断里只要距上次被接受的事件不到 20 ms 就丢弃，之后没有任何东西再上报一次电平。\file{/dev/key0} 可能收不到松开事件；sysfs 的 \code{value} 现在直接读 PB4，因此仍能显示真实电平。Linux gpio-keys 的做法是在中断里启动消抖定时器，到期时再读电平上报。
  \item \textbf{GPIOB 父 IRQ 目前只能被一个驱动占用。} \code{t113\_request\_irq()} 的一个 IRQ 槽只保存一个 handler，而 PB0～PB31 共用 IRQ 101。以后应增加 GPIO irqchip/bank handler，由它读取 STATUS 并分发每个 pin 的虚拟 IRQ。
  \item \textbf{硬件消抖很短。} \code{PB\_EINT\_DEB = 0x71} 表示 24 MHz HOSC、/128，采样周期约 5.3 微秒；主要消抖仍由驱动的 20 ms 软件窗口完成。
  \item probe 中途失败时，已经注册的字符设备号、devfs 名字和 IRQ 都没有撤销（现在也没有注销接口），重试 probe 会因为重复注册失败。
  \item sysfs 的 inode 号是路径哈希。两条不同路径哈希相同时，VFS 的 inode 缓存会把第二条当成第一条，类型和 fops 可能用错。现在的路径只有几十条，碰撞概率很低。
\end{enumerate}
